% Part: proof-theory
% Chapter: sequent-calculus
% Section: quantifiers

\documentclass[../../../include/open-logic-section]{subfiles}

\begin{document}

\olfileid{pt}{seq}{qua}

\olsection{البراهين المنتظمة والاستبدال}

تُدخل قواعد الكميات بعض التعقيدات بسبب «شروط المتغير المميَّز» التي تنطبق
على قاعدتي~\LeftR{\lexists} و\RightR{\lforall}. ويمكن معالجة معظم هذه
التعقيدات بضمان ألا يعاد استعمال~!!{constant}~$a$ في هذه الاستدلالات.
فلنضع التعريف الآتي.

\begin{defn}
يكون~!!^a{proof}~$\pi$ في~\Log{G1c} أو~\Log{G3c} \emph{منتظمًا} إذا تحقق
ما يأتي:
\begin{enumerate}
  \item لا يُستعمل أي متغير مميَّز~$a$ متغيرًا مميَّزًا لأكثر من استدلال
  واحد بقاعدة~\LeftR{\lexists} أو~\RightR{\lforall}؛ و
  \item إذا كان~$a$ المتغير المميَّز لاستدلال بقاعدة~\LeftR{\lexists}
  أو~\RightR{\lforall}، فإنه لا يقع في~$\pi$ إلا فوق ذلك الاستدلال.
\end{enumerate}
\end{defn}

\begin{lem}\ollabel{lem:subst-regular} إذا كان~$\pi$ منتظمًا وينتهي
بـ$\Gamma \Sequent \Delta$، وكان~$c$~!!a{constant} لا يُستعمل متغيرًا
مميَّزًا في~$\pi$، وكان~$t$ حدًّا لا يحتوي أي متغير مميَّز في~$\pi$،
فإن~$\Subst{\pi}{t}{c}$ الناتج من~$\pi$ باستبدال كل وقوع لـ$c$ بـ$t$
هو~!!{proof} منتظم. ومتتاليته النهائية هي
$\Subst{\Gamma}{t}{c} \Sequent \Subst{\Delta}{t}{c}$، حيث تدل
$\Subst{\Gamma}{t}{c}$ على متعددة المجموعات الناتجة من استبدال~$c$ بـ$t$
في كل وقوع لصيغة في~$\Gamma$، مع إبقاء جميع الوقوعات وتعددياتها.
\end{lem}

\begin{proof}
  بالاستقراء على~$\pheight{\pi}$. إذا كان~$\pheight{\pi}=0$، اقتصر
  $\pi$ على بديهية. وعندئذ يكون~$\Subst{\pi}{t}{c}$ بديهية أيضًا، لأن
  كل~!!{formula}~$!A(c)$ فيها تتحول إلى~$!A(t)$.

  إذا كان~$\pheight{\pi}>0$، فنميّز الحالات بحسب الاستدلال الأخير
  في~$\pi$. وحالات القواعد البنيوية (في~\Log{G1c}) والقواعد المنطقية
  لروابط القضايا سهلة ونتركها تمارين. وننظر في الحالات التي ينتهي فيها
  $\pi$ باستدلال كمي.
  \begin{enumerate}
    \item إذا كان الاستدلال الأخير هو~\RightR{\lforall}، كان~$\pi$ من
    الصورة
    \[
    \AxiomC{}
    \RightLabel{$\pi'$}
    \Deduce$\Gamma(c) \fCenter \Delta'(c), !A(a,c)$
    \RightLabel{\RightR{\lforall}}
    \UnaryInf$\Gamma(c) \fCenter \Delta'(c), \lforall[x][!A(x,c)]$
    \DisplayProof
    \]
    وبفرض الاستقراء، يكون~$\Subst{\pi'}{t}{c}$~!!{proof} منتظمًا
    لـ$\Gamma(t) \fCenter \Delta'(t), !A(a,t)$. ولما كان~$a$ متغيرًا
    مميَّزًا في~$\pi$، فإنه لا يقع في~$t$. ومن ثم يكون الاستدلال الأخير
    في~$\Subst{\pi}{t}{c}$،
    \[
    \AxiomC{}
    \RightLabel{$\Subst{\pi'}{t}{c}$}
    \Deduce$\Gamma(t) \fCenter \Delta'(t), !A(a,t)$
    \RightLabel{\RightR{\lforall}}
    \UnaryInf$\Gamma(t) \fCenter \Delta'(t), \lforall[x][!A(x,t)]$
    \DisplayProof
    \]
    صحيحًا: فالنتيجة لا تحتوي~$a$.
    \item إذا كان الاستدلال الأخير هو~\LeftR{\lforall} في~\Log{G3c}،
    كان~$\pi$ من الصورة
    \[
    \AxiomC{}
    \RightLabel{$\pi'$}
    \Deduce$!A(s(c),c), \lforall[x][!A(x,c)], \Gamma(c) \fCenter \Delta'(c)$
    \RightLabel{\LeftR{\lforall}}
    \UnaryInf$\lforall[x][!A(x,c)], \Gamma(c) \fCenter \Delta'(c)$
    \DisplayProof
    \]
    وبفرض الاستقراء، يكون~$\Subst{\pi'}{t}{c}$~!!{proof} منتظمًا
    لـ$!A(s(t),t), \lforall[x][!A(x,t)], \Gamma(t) \fCenter
    \Delta'(t)$. ويكون الاستدلال الأخير في~$\Subst{\pi}{t}{c}$،
    \[
    \AxiomC{}
    \RightLabel{$\Subst{\pi'}{t}{c}$}
    \Deduce$!A(s(t),t), \lforall[x][!A(x,t)], \Gamma(t) \fCenter \Delta'(t)$
    \RightLabel{\LeftR{\lforall}}
    \UnaryInf$\lforall[x][!A(x,t)], \Gamma(t) \fCenter \Delta'(t)$
    \DisplayProof
    \]
    صحيحًا. وحالة~\LeftR{\lforall} في~\Log{G1c} مشابهة، لكنها أبسط
    قليلًا.
    \item إذا كان الاستدلال الأخير~\RightR{\lexists} أو
    \LeftR{\lexists} فهذه الحالة تمرين.
  \end{enumerate}
  وفي كل حالة أضفنا متتالية إلى نهاية~!!{proof} منتظم (أو إلى
  !!{proof}s منتظمة في حالة القواعد ثنائية المقدمة).
  في كل حالة لا تتغير استدلالات المتغير المميَّز. وبما أن~$c$ لم يكن
  متغيرًا مميَّزًا، وأن~$t$ لا يحتوي أي متغير مميَّز في~$\pi$، فلا يُدخل
  الاستبدال متغيرًا مميَّزًا تحت الاستدلال الذي يخصه؛ كما تحافظ فروض
  الاستقراء على الانتظام. ومن ثم يكون~$\Subst{\pi}{t}{c}$ منتظمًا.
\end{proof}

\begin{prop}
إذا كان~$\pi$~!!a{proof} في~\Log{G1c} أو~\Log{G3c}، فهناك
!!{proof} منتظم~$\pi'$ له البنية نفسها (وبخاصة~!!{height} نفسه)،
ولا يختلف عن~$\pi$ إلا باستبدال المتغيرات المميَّزة.
\end{prop}

\begin{proof}
لأغراض هذا البرهان فقط، نسمي استدلال متغير مميَّز في~$\pi$
\emph{نظيفًا} إذا لم يكن متغيره المميَّز متغيرًا مميَّزًا لاستدلال آخر،
ولم يقع في~$\pi$ خارج~!!{proof} المباشر للاستدلال؛ ونسميه
\emph{متسخًا} فيما عدا ذلك. وليكن~$n$ عدد استدلالات المتغير المميَّز
المتسخة في~$\pi$. نبين بالاستقراء على~$n$ أنه يمكن تحويل~$\pi$ إلى
!!{proof} منتظم~$\pi'$ من النوع المطلوب. إذا كان~$n=0$، كانت جميع
استدلالات المتغير المميَّز في~$\pi$ نظيفة، ومن ثم كان~$\pi$ منتظمًا،
فلا شيء يلزم إثباته. وفيما عدا ذلك، نختار أعلى استدلال متسخ لمتغير
مميَّز، وليكن مثلًا استدلالًا بقاعدة~\RightR{\lforall}. ويكون
الـ!!{proof} الجزئي~$\pi_1$ المنتهي به من الصورة
    \[
    \AxiomC{}
    \RightLabel{$\pi_2$}
    \Deduce$\Gamma \fCenter \Delta, !A(c)$
    \RightLabel{\RightR{\lforall}}
    \UnaryInf$\Gamma \fCenter \Delta, \lforall[x][!A(x)]$
    \DisplayProof
    \]
ولأن~$c$ متغير مميَّز، فإنه لا يقع في~$\Gamma$ أو~$\Delta$ أو
$\lforall[x][!A(x)]$. ويكون البرهان~$\pi_2$ منتظمًا لأن كل استدلال
لمتغير مميَّز يقع فوق الاستدلال المختار نظيف. ولتكن~$a$~!!a{constant}
لا تقع في~$\pi$. وبحسب~\olref{lem:subst-regular}، يكون
$\Subst{\pi_2}{a}{c}$ برهانًا منتظمًا لـ$\Gamma \fCenter \Delta,
!A(a)$، ويمكن تطبيق~\RightR{\lforall} عليه. وإذا استبدلنا~$\pi_1$
في~$\pi$ بالبرهان~$\pi_1'$،
    \[
    \AxiomC{}
    \RightLabel{$\Subst{\pi_2}{a}{c}$}
    \Deduce$\Gamma \fCenter \Delta, !A(a)$
    \RightLabel{\RightR{\lforall}}
    \UnaryInf$\Gamma \fCenter \Delta, \lforall[x][!A(x)]$
    \DisplayProof
    \]
نكون قد استبدلنا متغيرًا مميَّزًا جديدًا~$a$ بالمتغير~$c$. وينتج عن ذلك
عدد أقل قطعًا من~$n$ من استدلالات المتغير المميَّز المتسخة، لأن إعادة
التسمية قد تنظف أكثر من استدلال واحد. فتتبع النتيجة من فرض الاستقراء.
وحالات قواعد المتغير المميَّز الأخرى مماثلة.
\end{proof}

لن نستعمل القضية المثبتة للتو استعمالًا صريحًا، لكننا سنفترض أن جميع
!!{proof}s التي ننظر فيها منتظمة؛ فإن لم تكن كذلك أمكننا دومًا استبدال
المتغيرات المميَّزة لجعلها منتظمة.

\end{document}
